% Add "answers" option to print the answers
\documentclass[addpoints]{exam}

\usepackage[utf8]{inputenc}
\usepackage{calc}
\usepackage[dvipsnames]{xcolor}
\usepackage{ifthen}
\usepackage{tcolorbox}

% \makebox[\textwidth]{Name:\enspace\hrulefill}
%
% \vspace{5mm}
%\makebox[\textwidth]{Instructor’s name:\enspace\hrulefill}

\definecolor{MyLightGray}{gray}{0.60}

\newcommand{\truefalse}
{		\begin{checkboxes}
			 \choice True
			 \choice False. Explanation:
		\end{checkboxes}
		\vspace{10pt}
}


% Check if 'answers' is an option
\newcommand{\ifanswers}[2]{%
  \ifthenelse{\boolean{exam@answers}}{#1}{#2}%
}



\newcommand{\answerspace}[1]
{
	\ifprintanswers
	\else
		\setlength{\mylength}{10pt}%
		% \vspace{#1\mylength}
		% Make the grey box
		\par\nobreak\noindent
		\makebox[\linewidth][l]{%
			% set inner padding to zero, so only the frame shows
			{\setlength{\fboxsep}{0pt}%
			% draw a color frame, white (or transparent) inside
			\fcolorbox{MyLightGray}{white}{%
			% the inner box has zero width content but fixed height
			\parbox[t][#1\mylength][c]{%
				% subtract twice the frame‐thickness from the full width
				\dimexpr\linewidth-2\fboxrule\relax
			}{\strut}%
			}%
			}%
		}%
		\par
		\vspace{3\mylength}

  	\fi
}
\newlength{\mylength}

\newcommand{\code}[1]{{\texttt{#1}}}


\usepackage{color}
\usepackage{lastpage}
\usepackage{verbatim}
\usepackage{fancyvrb}
\usepackage{graphicx}
\usepackage{float}
\usepackage{minted}
\usepackage{multicol}

\RecustomVerbatimCommand{\VerbatimInput}{VerbatimInput}{}

\newcommand{\Code}[1]{\mintinline{python}{#1}}

\renewcommand{\solutiontitle}{\noindent\textbf{Solution sketch: }}
\definecolor{SolutionColor}{rgb}{0.8,0.9,1}

\cfoot{Page \thepage\ of \pageref{LastPage}}
\rhead{Initials: \rule{20mm}{0.4pt}}


%%%%%%%%%%%%%%%%%%%%%%%%%%%%%%%%%%%%%%%%
\begin{document}
\fvset{numbers=left,numbersep=3pt}
\thispagestyle{empty}

\begin{center}
	{\LARGE CS 408 Fall 2026 -- Test Case Minimization Quiz \\}
\fbox{\fbox{\parbox{5.5in}{\centering
	\begin{itemize}
        \item Try to come up with a solution on your own first, but
          after that talk to your friends and neighbors.

        \item Full points will be rewarded for \textbf{any} (reasonable)
          attempt.
        \end{itemize}
	}}}
\end{center}

\vspace{5mm}
\makebox[\textwidth]{\textbf{Name (please print)}: \rule{100mm}{0.8pt} }

\vspace{10mm}
\makebox[\textwidth]{\textbf{PUID}: \rule{100mm}{0.8pt} \hspace{35mm}}

\begin{questions}

\question A fuzzer has handed us an input file of \textbf{eight lines}
  that makes the program crash. Our goal is to find a smaller subset
  of those lines that cause a crash.

\begin{parts}

\part Suppose the crash is caused by \textbf{line 6 alone}. List the
  candidates \Code{ddmin} tests, in order, and say how many tests it
  takes. Write the candidates as lists of line numbers, e.g.,
  \Code{[1,2,3]} %
  \answerspace{10} %
  \begin{solution} %
    Five tests.

    \begin{center}\small
      \begin{tabular}{llll}
        & candidate & & outcome \\ \hline
        1 & \Code{[1,2,3,4]} & subset     & passes \\
        2 & \Code{[5,6,7,8]} & subset     & \textbf{crashes} --- keep it \\
        3 & \Code{[5,6]}     & subset     & \textbf{crashes} --- keep it \\
        4 & \Code{[5]}       & subset     & passes \\
        5 & \Code{[6]}       & subset     & \textbf{crashes} --- keep it \\
      \end{tabular}
    \end{center}

    Only step 2 of the algorithm ever fires. This is plain binary
    search, and it costs about $\log_2 n$ tests.
  \end{solution}

\part Now suppose the crash needs to include \textbf{both} lines 2 and
  7 \textbf{together}. What are the first four candidates \Code{ddmin}
  tests in this scenario, ignoring duplicates? (%
  \answerspace{20} %
  \begin{solution} %
    The first round makes no progress at all:

    \begin{center}\small
      \begin{tabular}{llll}
        & candidate & & outcome \\ \hline
        1 & \Code{[1,2,3,4]} & subset     & passes \\
        2 & \Code{[5,6,7,8]} & subset     & passes \\
        3 & \Code{[3,4,5,6,7,8]} & complement & passes \\
        4 & \Code{[1,2,5,6,7,8]} & complement & fails \\
      \end{tabular}
    \end{center}

    The two culprits are in different halves, so cutting anywhere in the
    middle destroys the crash. A pure binary search would stop here and
    report all eight lines.

    \textbf{Step 4 rescues it:} double the granularity to $n = 4$ and try
    again. Now the complements remove only a quarter at a time, and
    \Code{[1,2,5,6,7,8]} (dropping lines 3 and 4) still crashes --- real
    progress. The search eventually reaches \Code{[2,7]}.

    It is much more expensive: \textbf{34 tests} rather than 5.
  \end{solution}

\newpage

\part \Code{ddmin} guarantees its answer is \textbf{1-minimal}:
  removing any single remaining line makes the crash go away. Suppose
  the program crashes if lines 1 and 2 are \textbf{both} present,
  \emph{or} if line 3 is present, and our starting input is
  \Code{[1,2,3,4]}.

  What does \Code{ddmin} return, and why? %
  \answerspace{20} %
  \begin{solution} %
    It returns \Code{[1,2]}.

    \begin{center}\small
      \begin{tabular}{llll}
        & candidate & & outcome \\ \hline
        1 & \Code{[1,2]} & subset & \textbf{crashes} --- keep it, discard line 3 \\
        2 & \Code{[1]}   & subset & passes \\
        3 & \Code{[2]}   & subset & passes \\
        4 & \Code{[2]}   & complement & passes \\
        5 & \Code{[1]}   & complement & passes \\
      \end{tabular}
    \end{center}

    \Code{[1,2]} really is 1-minimal: drop either line and the crash
    goes away. But \Code{[3]} is \emph{smaller} and also crashes --- and
    \Code{ddmin} threw it away in its very first step and never looked
    back.

    \textbf{1-minimal means locally irreducible, not smallest.} When an
    input has two independent causes, delta debugging finds one of them
    and loses the other.
  \end{solution}

\part Why does \Code{ddmin} only guarantee 1-minimality, as opposed to
  the smallest failing example? %
  \answerspace{16} %
  \begin{solution} %
    Because finding the true smallest subset means searching the
    powerset: $2^n$ candidates for $n$ lines. For the eight-line file in
    part (a) that is 256 tests --- annoying but possible. For a
    4{,}000-byte input from a fuzzer it is $2^{4000}$, which is not
    ``slow'' but ``never''.

    1-minimality is reachable in $O(n^2)$ tests in the worst case, and
    far fewer in practice --- five tests in part (a). It is the price of
    an algorithm that terminates.
  \end{solution}

\end{parts}

\end{questions}

\end{document}
